Let $\omega_1,\omega_2$ be two circles that intersect at $X$ and $Y$, and let $A$ be a point on $\omega_1$. Lines $AX$ and $AY$ intersect $\omega_2$ at $B$ and $C$, respectively, such that $AY=YC=20$ and $AX=25$. $M$ is the midpoint of $AB$ and $N$ is the midpoint of $AM$. Denote by $\omega$ the circle passing through $M$ tangent to segments $AB$ and $BC$. $D$ is the tangency point of $\omega$ to $\overline{BC}$, and $\overline{AD}$ intersects $\omega$ again at $E$. Suppose that the perpendicular bisector of $DE$ passes through the midpoint of $BN$. Then $AE^2$ can be written in the form $\dfrac{m}{n}$, where $m$ and $n$ are relatively prime positive integers. Find $m+n$.
